% !TEX program = xelatex
\documentclass[11pt, a4paper]{article}

\usepackage[british]{babel}
\usepackage[inner=3cm, outer=2.5cm, top=2.5cm, bottom=3cm]{geometry}
\usepackage{fontspec}
\usepackage{setspace}
\usepackage[version=4]{mhchem}
\usepackage{siunitx}
\usepackage{enumitem}
\usepackage{titlesec}
\usepackage{titling}
\usepackage{qrcode}
\usepackage[hidelinks]{hyperref}

\setstretch{1.2}
\sisetup{locale = US, range-units=single, product-units=single}

\title{Advanced 2D and 3D characterisation of clinker phases and hydrated cementitious binders by combining a variety of modern imaging and analytical techniques}
\author{Florian Kleiner}

\begin{document}

%%% Page 1: Title Page %%%
\begin{titlepage}
    \centering
    \vspace*{1cm}
    
    {\Large Summary of the Doctoral Thesis\par}
    \vspace{1.5cm}
    
    {\LARGE \textbf{\thetitle}\par}
    \vspace{2cm}
    
    {\large Author:\par}
    {\Large \theauthor\par}
    \vspace{1cm}
    
    {\large Target degree:\par}
    {\Large Doktoringenieur (Dr.-Ing.)\par}
    \vspace{1cm}
    
    { at the\par}
    {\large Faculty of Civil and Environmental Engineering,\par}
    {\large Bauhaus-Universität Weimar\par}
    \vspace{1cm}

    {\large Supervisor:\par}
    {\Large Professor Dr.-Ing. Horst-Michael Ludwig\par}
    \vspace{1cm}
    
    {\large Status of the Doctoral Candidate:\par}
    {\Large Internal\par}
    \vspace{3cm}
    
    \vfill
    
    {\large Weimar, \today\par}
\end{titlepage}

\newpage

\section*{Summary of the Doctoral Thesis}

\subsection*{I. Problem Statement and Objectives}

\begin{enumerate}[label=\textbf{\arabic*.}, leftmargin=1.5cm, labelsep=0.5cm, itemsep=0.5em]
    \item Cement hydration is a multi-scale process.
    While qualitative microstructural observations are widely obtained, statistically representative, quantitative data analysis is not always the norm in cement-based materials science.
    \item Conventional image-based characterisation techniques often reach their resolution limits regarding nano-scale porosity, or fail to provide the chemical sensitivity required to assign trace elements to specific mineral phases within the heterogeneous clinker matrix.
    \item The primary objective of this thesis is to overcome these limitations by establishing correlative microscopy workflows and automated evaluation tools that bridge the gap between high-resolution imaging and physical property determination.
    \item Consequently, this combination of techniques should yield more information than each method can provide individually.
\end{enumerate}

\subsection*{II. State of the Art}

\begin{enumerate}[label=\textbf{\arabic*.}, leftmargin=1.5cm, labelsep=0.5cm, itemsep=0.5em, resume]
    \item \textbf{Preparation Artefacts:} Traditional mechanical polishing often introduces topographic relief and resin smearing, which distorts the representation of the native pore structure and complicates the segmentation of hydrate phases in Scanning Electron Microscopy (SEM).
    \item \textbf{3D Imaging Constraints:} Focused Ion Beam (FIB) nano-tomography (nT) is typically limited to very small, non-representative volumes.
    Existing workflows struggle with artefacts such as charging and redeposition, hindering the creation of high-quality datasets for microstructural modelling.
    \item \textbf{Subjectivity in Data Evaluation:} Many quantitative measures in cement research, such as the thickness of hydrate layers or particle spacing in fresh pastes, rely on manual or semi-automated measurements that are prone to operator bias and low statistical throughput.
\end{enumerate}

\subsection*{III. Methodology}

\begin{enumerate}[label=\textbf{\arabic*.}, leftmargin=1.5cm, labelsep=0.5cm, itemsep=0.5em, resume]
    \item \textbf{ArBIB Preparation:} This thesis utilises Argon Broad Ion Beam (ArBIB) polishing (at room temperature and at \qty{-140}{\celsius}) and milling as essential preparation steps for artefact-reduced imaging and characterisation.
    This methodology allows for the visualisation of native nanopores within dense C-S-H without the interference of embedding media.
    \item \textbf{Multi-Modal Correlation:} A central focus of this research is the integration of multi-modal data.
    This is achieved by co-registering SEM-BSE image data and SEM-EDX phase maps with trace element distributions obtained using Laser Ablation - Inductively Coupled Plasma - Mass Spectrometry (LA-ICP-MS) and mechanical property maps obtained by Nanoindentation (NI).
    The combination of these data allows the direct assignment of chemical and mechanical properties to the phases of a hydrated cement or clinker.
    \item \textbf{Automated Python-based tools:} The research presents a newly developed algorithm for the automated measurement of Hydrate Layer Thickness (HLT) from large-scale BSE images, enabling the statistical analysis of hydration progress across hundreds of thousands of individual particles.
    This method can also be adapted to measure particle distances and can be combined with Electron Backscatter Diffraction (EBSD) measurements to investigate whether a dependency exists between crystal orientation and hydrate layer growth.
    \item \textbf{Lift-Out Technique for FIB-nT:} The implementation of the lift-out technique for FIB-nT significantly reduces charging artefacts and stabilises the specimen.
    \item \textbf{Laser Preparation for FIB-nT:} The use of femtosecond (fs) laser ablation enables the rapid preparation of representative volumes for subsequent FIB-nT analyses.
    Furthermore, it allows for the fast preparation of large sections for Interfacial Transition Zone (ITZ) analysis, bridging the gap between nano-scale and bulk analysis.
\end{enumerate}

\subsection*{IV. Principal Results}

\begin{enumerate}[label=\textbf{\arabic*.}, leftmargin=1.5cm, labelsep=0.5cm, itemsep=0.5em, resume]
    \item \textbf{ArBIB Polishing:} It was demonstrated that the preparation of dense inner C-S-H at \qty{-140}{\celsius} results in smaller observable pores.
    This indicates a minor damaging effect of ArBIB polishing at room temperature.
    \item \textbf{Trace Element Mapping:} It was proven that trace elements such as Ba, V, and Rb preferentially incorporate into belite (\ce{C2S}), whereas Ti and Mn accumulate in the interstitial phases.
    \item \textbf{Micromechanical Phase Assignment:} The combination of NI-EDX allowed for a direct, non-statistical assignment of mechanical properties to specific hydrate compositions (C-S-H vs. C-A-S-H), in contrast to the traditional statistical deconvolution methods which can oversimplify the phase complexity of the hydrated binder.
    \item \textbf{Statistical Representativeness:} A systematic area analysis showed that the Representative Elementary Area for phase composition measurements scales with the size of image features: While approx. \qty{0.3}{\milli\meter\squared} is sufficient for homogeneous sample sections of hydrated \ce{C3S}, specimens with large-scale features (e.g., air voids) required over \qty{1.4}{\milli\meter\squared} for statistically reliable results.
    \item \textbf{Hydrate Layer Growth and Phase Separation:} The HLT analysis quantified the growth of hydration products over time, confirming that \ce{C2S} hydrates significantly slower than \ce{C3S}.
    In multi-phase mixtures, the algorithm successfully distinguished between the phases, identifying distinct HLT peaks that reflect the individual hydration kinetics of \ce{C3S} and \ce{C2S}.
    \item \textbf{Hydration Kinetics and Crystal Orientation:} Correlation of the HLT measurements with the crystal orientation of \ce{C3S} indicates that C-S-H growth around \ce{C3S} particles is governed much more by local spatial availability and pore connectivity than by crystallographic orientation.
    \item \textbf{Quantification of Dispersing Effectiveness:} The use of cryo-FIB-SEM combined with automated inter-particle distance analysis was able to illustrate the dispersing effect of power ultrasound and polycarboxylate ethers in fresh pastes, providing an objective tool for rheological optimisations.
    \item \textbf{Porosity Measurements:} A multi-method comparison (LT-DSC, mercury intrusion porosimetry, 2D, and 3D imaging) highlighted the advantages and disadvantages of each method. The LT-DSC analyses revealed a clear correlation over time between the reduction of capillary porosity and the formation of gel porosity in \ce{C3S} and \ce{C2S}.
\end{enumerate}

\vfill
\begin{center}
    \qrcode[height=2cm]{https://github.com/kleinerELM/Dissertation/releases}\hspace*{.5cm}\url{https://github.com/kleinerELM/Dissertation/releases}
\end{center}

\end{document}